% Image credits for the photographs and AI illustrations of Book 3
% (University Biology, Year 1). Machine-readable license records live in
% images/book3/CREDITS.md; keep the two files in sync.
\chapter*{Beeldverantwoording}

% Under bidi=basic a Latin run is ONE directional box, so a long author or
% institution name cannot be broken at all; \allowbreak inside it does nothing.
% The Arabic edition reported this page as the book's only overfull boxes.
% \sloppy must be in force when the paragraph is BROKEN, hence the \par before
% \endgroup -- without it the group closes first and the tolerance is restored.
\begingroup\sloppy

De tekeningen in dit boek zijn origineel. De foto's worden gebruikt
onder hun respectieve vrije licenties, met dank aan hun makers:

\begin{itemize}
  \item \emph{Claude Bernard} (hoofdstuk 1): foto, negentiende eeuw,
        Wikimedia Commons, publiek domein.
  \item \emph{Kurkcellen} (hoofdstuk 5): plaat uit Robert Hookes
        \emph{Micrographia}, 1665, Wikimedia Commons, publiek domein.
  \item \emph{Antoni van Leeuwenhoek} (hoofdstuk 5): schilderij van
        Jan Verkolje, ca.~1680, Rijksmuseum via Wikimedia Commons,
        publiek domein.
  \item \emph{Escherichia coli} (hoofdstuk 5): rasterelektronen\-microscopische
        opname, NIAID, via Wikimedia Commons, publiek domein.
  \item \emph{Linus Pauling} (hoofdstuk 12): foto, Nobel Foundation,
        1962, via Wikimedia Commons, publiek domein.
  \item \emph{Melvin Calvin} (hoofdstuk 14): foto, Lawrence Berkeley
        Laboratory, via Wikimedia Commons, publiek domein.
  \item \emph{Hans Krebs} (hoofdstuk 15): foto, Nobel Foundation,
        1953, via Wikimedia Commons, publiek domein.
  \item \emph{Menselijk karyogram} (hoofdstuk 17): National Human
        Genome Research Institute, via Wikimedia Commons, publiek
        domein.
  \item \emph{Rode bloedcel, bloedplaatje en lymfocyt} (hoofdstuk 7):
        rasterelektronen\-microscopische opname, Electron Microscopy
        Facility, National Cancer Institute, via Wikimedia Commons,
        publiek domein.
  \item Hoofdstuk~\ref{ch:b1:classifying-biodiversity}: portret van Carl Linnaeus door Alexander Roslin, 1775, Nationalmuseum, Stockholm --- publiek domein (Wikimedia Commons).
  \item Hoofdstuk~\ref{ch:b1:classifying-biodiversity}: foto van Carl Woese door Don Hamerman --- CC~BY~3.0 (Wikimedia Commons).
  \item Hoofdstuk~\ref{ch:b1:classifying-biodiversity}: twee platen uit Ernst Haeckels \emph{Kunstformen der Natur} (1904), Diatomea en Cyrtoidea --- publiek domein (Wikimedia Commons).
\end{itemize}

De volledige bronvermeldingen en licentieadressen van elke foto staan
in \texttt{images/\allowbreak book3/\allowbreak CREDITS.md} in de repository van het boek.

\bigskip
Naast de foto's en de originele tekeningen zijn ongeveer zeventig
figuren verspreid over de hoofdstukken van dit deel illustraties die
met kunstmatige intelligentie zijn gemaakt, met de beeldgeneratie van
OpenAI en op basis van prompts die de redacteuren van het boek hebben
geschreven, en die vóór opname op biologische juistheid zijn
nagekeken. De volledige prompt van elk gegenereerd beeld staat in
\texttt{images/\allowbreak book3/\allowbreak ai/\allowbreak PROMPTS.md} in de repository.
\par\endgroup
\clearpage
